\subsection{Definición de los indicadores espectrales y de la prueba 0--1}
\label{bre:diag:definition}

Para una componente postransitoria $u_j$, muestreada uniformemente con
intervalo $\Delta t$, se sustrae su media $\bar u$ y se aplica una ventana
$w_j$. Con $N$ muestras, la transformada discreta es
\begin{equation}
 S_k=\sum_{j=0}^{N-1}w_j(u_j-\bar u)e^{-2\pi\ii jk/N},
 \qquad f_k=\frac{k}{N\Delta t}.
 \label{bre:diag:dft}
\end{equation}
La ventana de Hann es $w_j=[1-\cos(2\pi j/(N-1))]/2$.
La representación de módulo emplea $|S_k|/\sum_jw_j$; es distinta de
una densidad espectral de potencia. Para una PSD unilateral por Welch,
cada segmento usa el factor $\Delta t/\sum_jw_j^2$ sobre $|S_k|^2$;
se duplican las frecuencias positivas interiores, excepto Nyquist
cuando está presente, y se promedian los segmentos. Esta normalización
conserva las unidades de potencia por unidad de frecuencia.

Para cuantificar concentración espectral se excluye el modo cero y se
normalizan las potencias $P_k=|S_k|^2$ en los $m$ modos positivos:
\begin{equation}
 p_k=\frac{P_k}{\sum_{\ell=1}^mP_\ell},\qquad
 H=-\frac{\sum_{k=1}^mp_k\log p_k}{\log m}.
 \label{bre:diag:entropy}
\end{equation}
Se requiere $m>1$ y potencia total positiva; una señal constante se trata
por separado. Con $0\log0=0$, se obtiene $0\le H\le1$:
la cota inferior sigue de $p_k\le1$ y la superior de la concavidad del
logaritmo, que maximiza la entropía en $p_k=1/m$.
El cociente de potencia dominante es $\max_kp_k$, y la frecuencia
dominante corresponde al modo que lo alcanza. Para comparar ventanas se
conservan paso de muestreo, longitud y ventana; de otro modo cambian la
resolución $1/(N\Delta t)$ y el número de modos usados en $H$.
La deriva relativa usada en el clasificador compara los máximos de
ventanas parciales mediante $(f_{\max}-f_{\min})/f_{\max}$, cuando
$f_{\max}>0$ y todas las frecuencias parciales son positivas. Si alguna
es nula o no puede estimarse, el clasificador asigna deriva infinita y
no acepta el filtro de periodicidad. Una señal transitoria o una resolución insuficiente pueden
producir banda ancha sin demostrar caos.

La prueba 0--1 utiliza al menos $N=100$ muestras finitas de una señal
preparada $\phi_j$, después de
declarar el descarte, submuestreo, eliminación de tendencia y
normalización \cite{gottwald2009}. Para $c\in(0,\pi)$, define
\begin{align}
 p_c(n)&=\sum_{j=1}^n\phi_j\cos(jc),\nonumber\\
 q_c(n)&=\sum_{j=1}^n\phi_j\sin(jc).
 \label{bre:diag:translations}
\end{align}
El desplazamiento cuadrático medio a un retardo $n<N$ es
\begin{equation}
 M_c(n)=\frac1{N-n}\sum_{j=1}^{N-n}
 \bigl[\Delta_np_c(j)^2+\Delta_nq_c(j)^2\bigr],
 \label{bre:diag:msd}
\end{equation}
donde $\Delta_np_c(j)=p_c(j+n)-p_c(j)$, y análogamente para $q_c$.
Se elimina el término oscilatorio de la media mediante
\begin{equation}
 D_c(n)=M_c(n)-\bar\phi^{\,2}
                 \frac{1-\cos(nc)}{1-\cos c}.
 \label{bre:diag:modified-msd}
\end{equation}
La corrección procede de sumar la progresión geométrica
$\sum_{j=1}^n e^{\ii jc}$ y calcular su módulo al cuadrado:
$|\sum_{j=1}^n e^{\ii jc}|^2=(1-\cos nc)/(1-\cos c)$.
Por tanto, resta exactamente la contribución de una señal constante.

Para el vector de retardos $n=1,\ldots,n_{\max}$ se calcula
$K_c=\operatorname{corr}(n,D_c(n))$. El valor comunicado es la mediana
de $K_c$ sobre los valores de $c$ elegidos; se conserva esa selección al
comparar ensayos. En el procedimiento de correlación utilizado,
$n_{\max}=\min(\max(20,\lfloor N/10\rfloor),500)$.
Una varianza nula del desplazamiento no admite correlación de Pearson;
se trata como caso degenerado regular. Los umbrales operativos
$K<0.2$, $K>0.8$ y el intervalo intermedio describen, respectivamente,
compatibilidad regular, compatibilidad caótica e indeterminación de ese
diagnóstico finito. Un valor ligeramente negativo es posible por la
correlación muestral. Ruido, transitorios y sobremuestreo afectan el
resultado; la mediana y la comparación de varios intervalos de muestreo
reducen sensibilidades concretas, sin demostrar caos ni ocultedad.
